\scanpage{240}
\deleted{i.e.}
\[
\deleted{gs_0=s_0,\quad g\check s_0=\lambda\check s_0\quad\text{i.e. }\unsure{g}\check s_{0}=g^{-1}\lambda\check s_0=\lambda g^{-1}\check s_0}
\]

en effet, comme $G_0$ opère transitivement sur $A_0$ et $A_\lambda$ (il suffit \inserted{donc} pour cet argument que l'on ait cette transitivité \textquotedblleft{}faible\textquotedblright, au lieu de la transitivité sur $\overrightarrow{A_0}$ \deleted{\unsure{fibre sur}} i.e. sur $R$), si on \deleted{\unsure{voit}} $A_0\cap A_\lambda=\varnothing$, on en conclut $A_0=A_\lambda$ d'où $\lambda=\mathrm{id}$, cqfd.

Nous voulons maintenant expliciter \deleted{\unsure{les diverses}} \inserted{la} \deleted{\unsure{situations|donnée}} \inserted{donnée} d'un \deleted{\unsure{de}} multigraphe à homothéties homogène (i.e. satisfaisant (59)), en termes de données de théorie des groupes --- plus particulièrement dans le cas où on suppose de plus satisfaite la condition d'isomorphie (58).

Pour ceci, \deleted{dans le \unsure{premier}} je vais prendre les choses par un bout légèrement différent, en \deleted{\unsure{me}} \uline{partant} d'un cas \textquotedblleft{}épinglé\textquotedblright, \deleted{graphe} savoir

\[
\left\{\begin{aligned}
&\text{a) graphe }\deleted{\unsure{\gamma}}\ \gamma_0=(S,A_0),\text{ muni d'\deleted{une arête}}\\
&\text{b) groupe }\deleted{\unsure{G_0}}H\text{ d'autom de }(S,A_0),\text{ transitif sur }\overrightarrow{A_0}\\
&\quad\inserted{[\text{\uline{et} sur }S,\ \text{i.e. }\forall s\in S,\ \mathrm{Et}\ A_0(s)\ne\varnothing]}\\
&\text{c) et d'une \textquotedblleft{}origine\textquotedblright{} }(s_0,t_0),\text{ avec donc }\{s_0,t_0\}\in A_0\\
&\quad\text{et enfin}\\
&\text{d) d'un sous-groupe } Z\subset\unsure{\mathfrak S|\mathcal G}_S\text{ commutatif,}\\
&\quad\text{formé d'homothéties pour le graphe }\gamma_0\\
&\quad\text{commutant à l'action de }\deleted{\unsure{G_0}}H,
\end{aligned}\right.
\tag*{(62)}
\]

sans nous préoccuper pour le moment \unsure{si|où} \deleted{\unsure{$G_0$}}\inserted{$H$} est \inserted{le groupe $G_0$}, formé de \uline{tous} les autom. \deleted{qui commutent} \inserted{de la structure} qui commutent à $Z$ --- et en fait je serai amené à travailler avec la donnée d'un sous-groupe \deleted{\unsure{du}} plus petit, $G_0^{+}$. \deleted{Autant alors partir tout de suite d'un tel sous-groupe, je vais le noter $H$, pour identifier la donnée correspondante. On voit que}

% archive sheet number: 239
